% =============================================================================
% test-abnt-resumo.tex
% Documento de teste para abnt-resumo.sty e folha de aprovação
%
% Demonstra resumo (NBR 6028:2021) e folha de aprovação (NBR 14724:2024
% seção 4.2.1.3), integrados com abnt-trabalho-academico.cls.
%
% Compilação:
%   TEXINPUTS=./src: latexmk -pdf -output-directory=tests tests/test-abnt-resumo.tex
%
% Checklist:
%   1  — Título "RESUMO" centralizado, negrito, maiúsculas
%   2  — Texto do resumo em parágrafo único
%   3  — Palavras-chave precedidas de "Palavras-chave:"
%   4  — Palavras-chave separadas por ponto e vírgula, finalizadas por ponto
%   5  — Título "ABSTRACT" centralizado, negrito, maiúsculas
%   6  — Keywords precedidas de "Keywords:"
%   7  — Folha de aprovação com nome do autor e título
%   8  — Natureza do trabalho na folha de aprovação
%   9  — Data de aprovação
%   10 — Linhas de assinatura com nomes e titulações
%   11 — Placeholders visíveis quando membros não definidos
%   12 — Compatibilidade com abnt-trabalho-academico.cls
%   13 — Ordem correta dos elementos pré-textuais
% =============================================================================

\documentclass{abnt-trabalho-academico}
\usepackage{abnt-resumo}

% Metadados
\titulo{Impactos ambientais da urbanização em áreas de preservação permanente}
\subtitulo{um estudo de caso na região metropolitana}
\autor{João Carlos da Silva}
\orientador{Prof. Dr. Maria Helena de Souza}
\instituicao{Universidade de São Paulo}
\local{São Paulo}
\ano{2025}
\natureza{Dissertação apresentada ao Programa de Pós-Graduação em
Geografia da Universidade de São Paulo, como requisito parcial para
obtenção do título de Mestre em Geografia.}
\programa{Programa de Pós-Graduação em Geografia}
\area{Geografia Física}

% Banca examinadora (checklist 10)
\dataaprovacao{15 de março de 2025}
\membrodabanca{Prof. Dr. Maria Helena de Souza (Orientadora)}
  {Doutora em Geografia \textendash\ Universidade de São Paulo}
\membrodabanca{Prof. Dr. Carlos Eduardo Lima}
  {Doutor em Ciências Ambientais \textendash\ Universidade Estadual de Campinas}
\membrodabanca{Prof. Dr. Ana Paula Oliveira}
  {Doutora em Planejamento Urbano \textendash\ Universidade Federal do Rio de Janeiro}

\begin{document}

% =====================================================================
% PRÉ-TEXTUAIS (checklist 13 — ordem correta)
% =====================================================================
\pretextual

% Capa
\imprimircapa

% Folha de rosto
\imprimirfolhaderosto

% Folha de aprovação (checklist 7, 8, 9, 10)
\imprimirfolhadeaprovacao

% Dedicatória
\imprimirdedicatoria{%
  A todos os pesquisadores que dedicam suas carreiras à preservação
  do meio ambiente.
}

% Agradecimentos
\imprimiragradecimentos{%
  Agradeço à minha orientadora, Prof. Dr. Maria Helena de Souza, pela
  dedicação e paciência durante todo o processo de pesquisa.

  Aos membros da banca examinadora, pelas valiosas contribuições.

  À CAPES, pelo apoio financeiro que viabilizou esta pesquisa.
}

% Epígrafe
\imprimirepigrafe{%
  ``A Terra não é uma herança dos nossos pais, mas um empréstimo dos
  nossos filhos.'' \\[0.5cm]
  \textit{Provérbio indígena}
}

% Resumo em português (checklist 1, 2, 3, 4)
\imprimirresumo{%
  Esta dissertação investiga os impactos ambientais decorrentes do processo
  de urbanização em áreas de preservação permanente na região metropolitana
  de São Paulo. A pesquisa utiliza uma abordagem metodológica mista,
  combinando análise de imagens de sensoriamento remoto, levantamento de
  campo e entrevistas com moradores das áreas afetadas. Os resultados
  indicam que o avanço da ocupação urbana irregular provocou degradação
  significativa da cobertura vegetal, alteração do regime hídrico e aumento
  da vulnerabilidade a eventos climáticos extremos. A análise temporal
  revela que a taxa de desmatamento acelerou nos últimos quinze anos,
  coincidindo com o crescimento populacional nas áreas periféricas. O estudo
  propõe um conjunto de diretrizes para compatibilizar a ocupação urbana com
  a preservação ambiental, incluindo mecanismos de regularização fundiária
  sustentável e programas de recuperação de áreas degradadas.
}{gestão ambiental; urbanização; áreas de preservação permanente;
  sensoriamento remoto; região metropolitana de São Paulo}

% Abstract em inglês (checklist 5, 6)
\imprimirabstract{%
  This dissertation investigates the environmental impacts of urbanization
  in permanent preservation areas in the metropolitan region of São Paulo.
  The research uses a mixed methodological approach, combining remote
  sensing image analysis, field surveys, and interviews with residents of
  affected areas. The results indicate that the advance of irregular urban
  occupation caused significant degradation of vegetation cover, alteration
  of the water regime, and increased vulnerability to extreme weather
  events. The temporal analysis reveals that the deforestation rate
  accelerated over the past fifteen years, coinciding with population
  growth in peripheral areas. The study proposes a set of guidelines for
  reconciling urban occupation with environmental preservation, including
  sustainable land regularization mechanisms and degraded area recovery
  programs.
}{environmental management; urbanization; permanent preservation areas;
  remote sensing; São Paulo metropolitan region}

% Sumário
\imprimirsumario

% =====================================================================
% ELEMENTOS TEXTUAIS
% =====================================================================
\textual

\chapter{Introdução}

Este trabalho aborda os impactos ambientais da urbanização.

\chapter{Revisão da literatura}

A literatura sobre o tema é vasta.

\chapter{Metodologia}

A pesquisa adotou abordagem mista.

\chapter{Resultados e discussão}

Os resultados confirmam a hipótese inicial.

\chapter{Considerações finais}

O estudo contribui para o debate sobre gestão ambiental urbana.

% =====================================================================
% PÓS-TEXTUAIS
% =====================================================================
\postextual

\chapter*{Referências}
\addcontentsline{toc}{chapter}{REFERÊNCIAS}

ASSOCIAÇÃO BRASILEIRA DE NORMAS TÉCNICAS. \textbf{NBR 6028}: informação e
documentação: resumo, resenha e recensão: apresentação. Rio de Janeiro, 2021.

\end{document}
